\documentclass[11pt,a4paper]{article}

\usepackage[utf8]{inputenc}
\usepackage[T1]{fontenc}
\usepackage{lmodern}
\usepackage{amsmath,amssymb,amsthm}
\usepackage{bm}
\usepackage[margin=2.5cm]{geometry}
\usepackage{hyperref}

\newcommand{\RR}{\mathbb{R}}
\newcommand{\CC}{\mathbb{C}}
\newcommand{\ZZ}{\mathbb{Z}}
\newcommand{\Tr}{\operatorname{Tr}}
\newcommand{\curl}{\operatorname{curl}}
\renewcommand{\div}{\operatorname{div}}

\title{Technical Research Note\\[4pt]
\large A Proposed Mathematical Bridge Between Hopf--Beltrami Geometry\\
and Quantum Morphodynamics}
\author{Fabian Leo Naressi\\
\small Independent Researcher, Senemosìa Punto Zero, Italy\\
\small Research framework: Quantum Morphodynamics (Morfodinamica Quantistica)}
\date{10 October 2026}

\begin{document}
\maketitle

\begin{abstract}
We outline a research programme linking the contact geometry of Hopf fibrations,
Beltrami fields on $S^3$, torsion deformations of a connection, and the spectral
structure of the discrete PRISM state space. For each step we state which
relations are exact consequences of the definitions, which are conjectural, and
which operator specifications are still missing. Three initial tests (A, B, C)
are proposed to decide which correspondences are exact, which are approximate
and which fail. No gravitational interpretation is claimed at this stage.
\end{abstract}

\section{Hopf fibration and contact geometry}

The finite-dimensional complex Hopf fibrations are
\begin{equation}
S^1\hookrightarrow S^{2n+1}\xrightarrow{\ \pi\ }\CC\mathrm{P}^{n},
\end{equation}
with the case $n=1$
\begin{equation}
S^1\hookrightarrow S^3\longrightarrow \CC\mathrm{P}^{1}\cong S^2
\end{equation}
and the higher-dimensional case
\begin{equation}
S^1\hookrightarrow S^9\longrightarrow \CC\mathrm{P}^{4}.
\end{equation}

Writing $z=(z_0,\ldots,z_n)\in\CC^{n+1}$ with $\sum_j|z_j|^2=1$, a standard
contact form on $S^{2n+1}$ is
\begin{equation}
\alpha=\frac{1}{2i}\sum_{j=0}^{n}\bigl(\bar z_j\,dz_j-z_j\,d\bar z_j\bigr).
\end{equation}
The contact distribution is $\xi=\ker\alpha$, and the contact condition is
\begin{equation}
\alpha\wedge(d\alpha)^n\neq0 .
\end{equation}
The Reeb vector field $R$ is characterized by
\begin{equation}
\alpha(R)=1,\qquad \iota_R\,d\alpha=0 .
\end{equation}
These equations provide the geometric starting point for comparing fiber
geometry, admissible modes and dynamics.

\section{Beltrami fields and geometric mode selection}

On an oriented Riemannian three-manifold, a Beltrami field satisfies
\begin{equation}
\curl B=\lambda B,
\end{equation}
or, in differential forms,
\begin{equation}
{*}\,dB^\flat=\lambda B^\flat ,
\end{equation}
where $B^\flat$ is the one-form metrically dual to $B$ and $\lambda$ is a
spectral parameter. Taking the divergence,
\begin{equation}
\div(\curl B)=0\quad\Longrightarrow\quad\lambda\,\div B=0,
\end{equation}
so that, for $\lambda\neq0$,
\begin{equation}
\div B=0 .
\end{equation}

For an incompressible fluid of constant density, the Euler equation is
\begin{equation}
\partial_tB+(B\cdot\nabla)B=-\nabla p,\qquad \nabla\cdot B=0,
\end{equation}
with $B$ interpreted as the velocity field. The vector identity
\begin{equation}
(B\cdot\nabla)B=\nabla\!\left(\frac{|B|^2}{2}\right)-B\times(\nabla\times B)
\end{equation}
shows that a Beltrami field satisfies
\begin{equation}
B\times(\nabla\times B)=0 .
\end{equation}
Consequently its nonlinear inertial term is a gradient and can be balanced by
the pressure in a stationary solution.

This gives a precise fluid-dynamical setting in which geometric constraints can
select special modes. On the round $S^3$, the relationship between Hopf fibers
and Beltrami fields can be investigated through the curl spectrum.

\section{Connection, curvature and torsion}

Let $P\to M$ be a principal bundle with connection one-form $A$. Its curvature
is
\begin{equation}
F=dA+A\wedge A .
\end{equation}
Under a gauge transformation $g$,
\begin{equation}
A\mapsto g^{-1}Ag+g^{-1}dg,\qquad F\mapsto g^{-1}Fg .
\end{equation}

For a geometric connection with coframe $e^a$ and spin connection
$\omega^a{}_b$, the torsion two-form is
\begin{equation}
T^a=De^a=de^a+\omega^a{}_b\wedge e^b,
\end{equation}
and the curvature two-form is
\begin{equation}
R^a{}_b=d\omega^a{}_b+\omega^a{}_c\wedge\omega^c{}_b .
\end{equation}
The torsion-free condition is $T^a=0$. A possible research deformation is
\begin{equation}
\omega^a{}_b(\varepsilon)=\omega^a{}_b(0)+\varepsilon K^a{}_b,
\end{equation}
where $K^a{}_b$ is a contorsion one-form and $\varepsilon$ a deformation
parameter.

For a suitably defined self-adjoint spectral operator
\begin{equation}
\mathcal B_\varepsilon=\mathcal B_0+\varepsilon V,
\end{equation}
an isolated eigenvalue $\lambda_k(0)$ has, under the usual perturbation
assumptions, first-order variation
\begin{equation}
\lambda_k'(0)=\langle u_k,Vu_k\rangle
\end{equation}
for a normalized, nondegenerate eigenvector $u_k$. For a degenerate eigenspace
one must instead diagonalize the restriction of $V$ to that eigenspace. This is
the appropriate mathematical test for the proposed connection between torsion
and spectral splitting.

\paragraph{Operator specification.}
On a three-dimensional manifold, ${*}d$ maps one-forms to one-forms. In nine
dimensions the ordinary Hodge operator instead maps a two-form to a seven-form.
Therefore a generalized Beltrami operator on the contact distribution in $S^9$
must be defined with its precise domain, codomain, metric and contact (or
horizontal) projection. This definition is necessary before asserting
self-adjointness, ellipticity or a discrete spectrum.

\section{PRISM state space and graph Laplacian}

The PRISM state space uses the fixed coordinate order
\begin{equation}
(G,L,C,T,B)\in\ZZ_{64}\times\ZZ_6^3\times\ZZ_5,
\end{equation}
with cardinality
\begin{equation}
N=64\cdot6\cdot6\cdot6\cdot5=69120 .
\end{equation}
The linear indexing convention is
\begin{equation}
\operatorname{ind}(G,L,C,T,B)=((((G\cdot6+L)\cdot6+C)\cdot6+T)\cdot5+B).
\end{equation}
The product graph is
\begin{equation}
\mathcal G=C_{64}\,\square\,C_6\,\square\,C_6\,\square\,C_6\,\square\,C_5,
\end{equation}
with combinatorial Laplacian
\begin{equation}
L_{\mathcal G}=L_{64}\otimes I\otimes I\otimes I\otimes I
+I\otimes L_6\otimes I\otimes I\otimes I+\cdots
+I\otimes I\otimes I\otimes I\otimes L_5 .
\end{equation}
For a cycle $C_m$ the eigenvalues are
\begin{equation}
\lambda_k(C_m)=4\sin^2\!\left(\frac{\pi k}{m}\right),\qquad k=0,\ldots,m-1,
\end{equation}
hence the product-graph spectrum is
\begin{equation}
\lambda_{g,l,c,t,b}=4\sin^2\!\frac{\pi g}{64}+4\sin^2\!\frac{\pi l}{6}
+4\sin^2\!\frac{\pi c}{6}+4\sin^2\!\frac{\pi t}{6}+4\sin^2\!\frac{\pi b}{5}.
\end{equation}
The graph is $10$-regular and its Laplacian is positive semidefinite.

The PRISM shift is defined by the orbit--phase representation
\begin{equation}
\Omega\simeq\ZZ_{108}\times\ZZ_{640},\qquad
F(o,p)=\bigl(o,\;p+1 \bmod 640\bigr).
\end{equation}
It follows directly that $F^{640}=I$, with $108$ cycles of length $640$
(indeed $108\cdot640=69120=N$). This is a discrete dynamical structure; a
physical interpretation requires an additional map connecting graph states to
physical fields.

\section{Spectral operator and the proposed bridge}

Let $\mathcal H$ be a Hilbert space of fields on the chosen geometric domain.
A candidate operator combining graph and geometric contributions is
\begin{equation}
L_{\mathrm{tot}}(\varepsilon)=a\,L_{\mathcal G}+b\,L_{\mathrm{geom}}
+\varepsilon\,V_{\mathrm{tors}},\qquad a,b\geq0,
\end{equation}
where $V_{\mathrm{tors}}$ is a specified torsion-dependent operator. The
coefficients and domains must be chosen so that all terms act on the same space
and have compatible units. If $L_{\mathrm{tot}}$ is self-adjoint and bounded
below, its spectral problem is
\begin{equation}
L_{\mathrm{tot}}(\varepsilon)\psi_k=\Lambda_k(\varepsilon)\psi_k .
\end{equation}

The key proposed question is whether graph modes can be mapped into geometric
modes through an explicitly defined map
\begin{equation}
\Phi:\mathcal H_{\mathrm{PRISM}}\longrightarrow\mathcal H_{\mathrm{geom}},
\end{equation}
with a controlled relation such as $L_{\mathrm{geom}}\Phi\approx\Phi L_{\mathcal G}$.
An exact intertwining relation would be
\begin{equation}
L_{\mathrm{geom}}\Phi=\Phi L_{\mathcal G}.
\end{equation}
This would imply that the image under $\Phi$ of a graph eigenvector is a
geometric eigenvector with the same eigenvalue, provided the image is nonzero
and the operator domains are respected. Whether such a map exists is an open
mathematical question for this proposed construction.

\section{Silver-Mean scaling and trace-map tests}

The silver ratio is $\delta_S=1+\sqrt2$, and its inverse-square scaling factor is
\begin{equation}
\delta_S^{-2}=(1+\sqrt2)^{-2}=3-2\sqrt2\approx0.1715728753 .
\end{equation}
For a sequence of spectral band widths $w_k$, define $r_k=w_{k+1}/w_k$. A
Silver-Mean scaling hypothesis would require evidence that
\begin{equation}
\lim_{k\to\infty}r_k=\delta_S^{-2},
\end{equation}
or a different explicitly justified scaling law. A few close numerical ratios do
not establish this limit.

For a transfer matrix $M_k(E)\in SL(2,\RR)$ define its trace
$x_k(E)=\Tr M_k(E)$. Since $\det M_k=1$, its eigenvalues $\mu,\mu^{-1}$ obey
\begin{equation}
\mu+\mu^{-1}=x_k .
\end{equation}
For a specified substitution, the exact matrix concatenation rule must first be
established. Trace recurrences then follow from the corresponding matrix
identities; they should not be assumed independently of the substitution
convention.

If a candidate trace map is
\begin{equation}
F(x,y,z)=(2x^2y-z,\;x,\;y),
\end{equation}
its Jacobian is
\begin{equation}
DF(x,y,z)=
\begin{pmatrix}
4xy & 2x^2 & -1\\
1 & 0 & 0\\
0 & 1 & 0
\end{pmatrix},\qquad \det DF=-1 .
\end{equation}
Its fixed points satisfy $x=y=z$ and $2x^3-2x=0$, so that
\begin{equation}
(x,y,z)\in\{(0,0,0),\,(1,1,1),\,(-1,-1,-1)\}.
\end{equation}
This map is only a candidate until it has been derived from the precise
transfer-matrix recursion used in the spectral model.

The numerical audit should track the same band through successive generations,
record its parent--child identity, and compare the measured widths with the
predicted scaling factor. Control sequences and a preregistered branch-selection
rule are necessary to rule out accidental matches.

\section{Emergent metric: a testable proposal}

Let $L(q)$ be a positive, invertible, self-adjoint operator depending smoothly
on geometric coordinates $q^\mu$. Define
\begin{equation}
A_\mu=L^{-1}\partial_\mu L .
\end{equation}
A candidate information-geometric tensor is
\begin{equation}
g_{\mu\nu}=\operatorname{Re}\Tr\bigl(A_\mu^\dagger A_\nu\bigr).
\end{equation}
For any real vector $v^\mu$,
\begin{equation}
v^\mu g_{\mu\nu}v^\nu=\Tr\!\left[(v^\mu A_\mu)^\dagger(v^\nu A_\nu)\right]\geq0,
\end{equation}
whenever the trace is well defined. This establishes positive
semidefiniteness under the stated assumptions, not nondegeneracy or a Lorentzian
signature.

A regularized version may be needed for infinite-dimensional operators; one
possibility is a heat-kernel or spectral-zeta prescription, provided the
relevant analytic conditions hold. The candidate effective action is
\begin{equation}
\Gamma[L]=\frac12\Tr\log\!\left(\frac{L}{\mu^2}\right)+V_{\mathrm{eff}},
\end{equation}
with sign and normalization to be fixed by the chosen physical convention.
Variation yields, schematically,
\begin{equation}
\delta\Gamma=\frac12\Tr\!\left(L^{-1}\delta L\right)+\delta V_{\mathrm{eff}},
\end{equation}
whenever the regularized trace and variation are justified.

A gravitational interpretation would require deriving the relevant field
equations, identifying the signature and causal structure, and showing how an
Einstein--Hilbert term emerges. A positive information metric alone is not a
derivation of general relativity.

\section{Proposed collaboration tests}

We suggest three initial tests.

\begin{description}
\item[Test A --- Geometric operator.] Specify the Hilbert space and the
generalized Beltrami operator on the relevant Hopf shell; prove or numerically
verify its symmetry, domain properties and spectral stability.
\item[Test B --- Torsion response.] Define $V_{\mathrm{tors}}$ from an explicit
connection deformation and compare perturbative eigenvalue shifts with direct
numerical diagonalization.
\item[Test C --- PRISM correspondence.] Construct and test an explicit map
$\Phi$ between PRISM graph modes and geometric modes. Separately test whether
the same branch of spectral bands exhibits the predicted Silver-Mean scaling.
\end{description}

The purpose is to identify which correspondences are exact, which are
approximations and which fail. A successful result would be a reproducible
mathematical bridge between bundle geometry, fluid-like dynamics and spectral
organization, with any gravitational interpretation developed only after the
operator and geometric steps have been established.

\end{document}
